\scanpage{45}
% archive sheet number: 44
\[
\mathbb{D}^{\handwritten{*}}(A)_{S'}=\deleted{R^1 f}\underline{H}_{\mathrm{DR}}^1(A'/S')\quad\deleted{\unsure{,|;}}.
\]
\deleted{\unsure{o|6}}On notera que d'après un théorème connu (cf. li\deleted{\unsure{b}}vre de Mumford Geometric Invariant theory) il existe toujours un $A'$ pour $S$ affine; le fait que \deleted{\unsure{il|y}}\inserted{$H^1(A'/S')$} ne dépende, à isomorphisme près, pas du prolongement choisi de $A$, peut être considéré ici comme conséquence du tapis de la cohomologie cristalline.

En fait, utilisant la filtration de \deleted{De Rham}\inserted{Hodge} sur la cohomologie de De Rham, on trouve une filtration sur $\mathbb{D}^{\handwritten{*}}(A)_S$, savoir
\[
0\handwritten{\longrightarrow}R^{\unsure{o|0}} f_{\handwritten{*}}(\underline{\handwritten{\Omega}}^1_{A/S})\handwritten{\longrightarrow}\underline{H}_{\mathrm{DR}}^1(A/S)\handwritten{\longrightarrow}R^1 f_*(\underline{O}_S)\handwritten{\longrightarrow}0\quad,
\]
qu'on peut interpréter comme
\diagram{\begin{tikzcd}[ampersand replacement=\&,column sep=small,row sep=small]
(**)\quad 0\arrow[r] \& \underline{t}_A^{\handwritten{\vee}}\arrow[r]\arrow[d,equal] \& \mathbb{D}^{\handwritten{*}}(A)_S\arrow[r] \& \underline{t}_{A^*}\arrow[r] \&0\ \unsure{-|{}}\ .\\
\& \handwritten{\underline{\omega}_A}
\end{tikzcd}}

Le foncteur $(*)$ et la suite exacte $(**)$ sont fonctoriels en $A$ et compatib\deleted{\unsure{l}}les avec image inverse. On peut \deleted{envisager} considérer que le foncteur $(*)$ se factorise en
\[
\mathbb{D}^*:\quad\text{Schémab}(S)\handwritten{\longrightarrow}\mathrm{Crisloclib}_{\mathrm{nilp}}\mathrm{fil}(S)
\]
vers les c\deleted{\unsure{r}}ristaux en modules localement libres ``filtrés'' (en deux crans localement libres). Utilisant cette factorisation, on trouve que pour un objet $S'$ du site cristallin nilpotent de $S$, la\unsure{,|{}} donnée d'un prolongement de $A$ en un schéma \inserted{abélien} $A'$ sur $S'$ revient à la donnée d'un prolongement de la filtration c\deleted{\unsure{e}}anonique $(**\handwritten{)}$ de $D^*(A)_S$ en une filtration (à quotients localement libres\handwritten{)} de $D^*(A)_{S'}$. Plus précisément, on a le

\uline{Théorème de déformation pour les schémas abéliens}. Soient $\deleted{\unsure{E}}S$ un sch\deleted{\unsure{e}}éma, $S'$ un voisinage à puissances div. loc. nilp. de $S$, alors le foncteur naturel
\[
\begin{aligned}
\text{Schémab}(S')\handwritten{\longrightarrow}{}&
\text{catégorie des couples }(A,\deleted{\unsure{D}}\deleted{\unsure{E}}\mathrm{Fil}^1)\text{ d'un schéma abélien}\\
& A\text{ sur }S\text{ et d'un sous-module localement facteur direct}\\
&\mathrm{Fil}^1\ \inserted{\text{de }\mathbb{D}^*(A)_{S'}}\text{ prolongeant }\mathrm{Fil}^1\mathbb{D}^*(A)_S\handwritten{\simeq\underline{\omega}}_A
\end{aligned}
\]
est une équivalence de catégories.

\handwritten{b)} On peut donner une deuxiè\deleted{\unsure{p}}me interprétation du foncteur $\mathbb{D}^*$ sur les
